%! Author = markt
%! Date = 3/10/2025

% Preamble
\documentclass[11pt]{article}

% Packages
\usepackage{latexsym,amsfonts,amssymb,amsthm,amsmath, yfonts, cancel, listings, xcolor, graphicx}
\usepackage{color}
\usepackage{listings}

\author{Mark Stanley}

\definecolor{codegreen}{rgb}{0,0.6,0}
\definecolor{codegray}{rgb}{0.5,0.5,0.5}
\definecolor{codepurple}{rgb}{0.58,0,0.82}
\definecolor{backcolour}{rgb}{0.95,0.95,0.92}
\lstdefinestyle{mystyle}{
    backgroundcolor=\color{backcolour},
    commentstyle=\color{codegreen},
    keywordstyle=\color{magenta},
    numberstyle=\tiny\color{codegray},
    stringstyle=\color{codepurple},
    basicstyle=\ttfamily\footnotesize,
    breakatwhitespace=false,
    breaklines=true,
    captionpos=b,
    keepspaces=true,
    numbers=left,
    numbersep=5pt,
    showspaces=false,
    showstringspaces=false,
    showtabs=false,
    tabsize=2
}

\title{Math 443 Session 25}

% Document
\begin{document}
    \lstset{style=mystyle}

    \maketitle

    \subsection{Elementary First Order Scalar ODE}
    All first order scalar ODEs can be written in the form, where the actual ODE is the original equation $u(t)$.

    \subsection{Equation}
    \[\frac{du}{dt}=au\]\\
    Where $a \in \mathbb{R}$, and $u(t)$ is an unknown scalar function.

    \subsection{Solution}
    The general solution to this equation is
    \[u(t)=ce^{at}\]\\
    Where $c$ is the integration constant, and the equation is derived using separation of variables.
    The unique solution is given by the initial condition $u(0)=u_0$.
    \[u(t)=be^{a(t-t_0)}\]

    \subsection{First Order Dynamic System}
    A first-order dynamical system consists of a system of $n$ couple first-order ODEs
    \[\frac{du_1}{dt}=f_1(t, u_1, \ldots, u_n), \ldots, \frac{du_n}{dt}=f_n(t, u_1, \ldots, u_n)\]
    The system can be written in vector form as
    \[\frac{du}{dt}=f(t, u)\]
    Where $u$ is a vector of the form $u=(u_1, \ldots, u_n)^T$ and $f$ is a vector function of the form $f=(f_1, \ldots, f_n)^T$.

    \subsection{Eigenvalues and Eigenvectors}
    \textbf{Definition:} Let $A$ be an $n \times n$ matrix.
    A scalar $\lambda$ is called an eigenvalue of $A$ if there exists a nonzero vector $v \neq 0$, called an eigenvectpr, such that:
    \[Av = \lambda v\]
    This equation can be rewritten as:
    \[(A-\lambda I)v=0\]
    Where $I$ is the identity matrix.
    A nontrivial solution exists iff $A-\lambda I$ is singular, i.e. $\det(A-\lambda I)=0$.

    \subsection{Computational Approach}
    To compute eigenvalues and eigenvectors manually, we can use the following steps:
    \begin{enumerate}
        \item Solve the following $\det(A-\lambda I)=0$ for each eigenvalue.
        \item For each eigenvalue, compute the eigenvector by solving $(A-\lambda I)v=0$.
    \end{enumerate}

    \subsection{Eigenspaces and Eigenvalues}
    \textbf{Definition:} Given a real eigenvalue $\lambda \in \mathbb{R}$, the corresponding eigenspace $V_\lambda \subset \mathbb{R}^n$ is the subspace spanned by all eigenvectors associated with $\lambda$.
    Equivalently:
    \[V_\lambda=\{v \in \mathbb{R}^n: Av=\lambda v\}=\ker(A-\lmabda I)\]

    \subsection{Properties:}
    \begin{enumerate}
        \item $\lambda$ is an eigenvalue $\iff \ker(A-\lambda I)\neq {0}$
        \item Every nonzero element of $V_\lambda$ is an eigenvector associated with $\lambda$.
        \item A basis of the eigenspace $V_\lambda$ is the most economical way to represent it.
    \end{enumerate}

    \subsection{Eigenvalue Problems: Key Properties}
    Let $A\in \mathbb{R}^{n\times n}$, and consider the eigenvalue problem:
    \[Av=\lambda v, \quad v \neq 0\]
    \mathbb{Main Properties:}
    \begin{enumerate}
        \item Eigenvalues $\lambda$ can be real or complex, even if $A$ is real
        \item The set of all eigenvectors corresponding to a given eigenvalue $\lambda$ (along with 0) forms a subspace called the eigenspace.
        \item A matrix $A$ is singular iff $\lambda = 0$ is an eigenvalue.
        \item For real matrices, complex eigenvalues come in conjugate pairs.
        \item The eigenvalues of a matrix are the roots of its characteristic polynomial: \[\det(A-\lambda I)=0\]
        \item If $A$ is symmetric ($A=A^T$), then:
        \begin{enumerate}
            \item All eigenvalues are real
            \item Eigenvectors corresponding to distinct eigenvalues are orthogonal
        \end{enumerate}
    \end{enumerate}
\end{document}